\section*{Capítulo \ref{ch:b2:oceanic-islands} --- Islas oceánicas}

\begin{solution}{exo:b2:oceanic-islands:1}
(a) Etapa de preescudo, como Loihi hoy. (b) \omterm{def:b2:oceanic-islands:stages}{Etapa de escudo}. (c) Etapa de postescudo. (d) Etapa de
rejuvenecimiento. (e) Ya ninguna etapa eruptiva: el volcán se ha hundido bajo su arrecife, que es ahora un \omterm{def:b2:oceanic-islands:reefs}{atolón}.
\end{solution}

\begin{solution}{exo:b2:oceanic-islands:2}
(a) $2500 + 350\sqrt9 = \qty{3550}{m}$. (b) $350\times(\sqrt{25} - \sqrt9) = 350\times2 =
\qty{700}{m}$. (c) $350/(2\times5) = \qty{35}{m}$ por millón de años.
\end{solution}

\begin{solution}{exo:b2:oceanic-islands:3}
(a) $5000/17 = \qty{294}{kyr}$ alrededor de las islas principales; $30\,000/70 = \qty{429}{kyr}$ a lo largo de la
cordillera. (b) Hay deslizamientos que pasan inadvertidos: un volcán más joven puede enterrar la cicatriz y los derrubios de otro más antiguo,
y los depósitos antiguos quedan cubiertos de sedimento. El número verdadero es mayor y el intervalo verdadero,
más corto.
\end{solution}

\begin{solution}{exo:b2:oceanic-islands:4}
(a) $4158\times0.3048 = \qty{1267}{m}$; $4610\times0.3048 = \qty{1405}{m}$; $150\times0.3048 =
\qty{46}{m}$. (b) El primer arrecife no creció a más de \qty{46}{m} de profundidad y su base está hoy a
\qty{1267}{m}: el volcán se ha hundido al menos $1267 - 46 = \qty{1221}{m}$ (más si el nivel del mar
era entonces más alto). (c) El hundimiento es como mucho de unos \qty{1267}{m}, en al menos \qty{33.9}{Myr}:
como mucho $1267/33.9 = \qty{37}{m}$ por millón de años, \qty{0.037}{mm} al año.
\end{solution}

\begin{solution}{exo:b2:oceanic-islands:5}
(a) $\dot S = 175/\sqrt t$: \qty{87.5}{m.Myr^{-1}} (\qty{0.0875}{mm.yr^{-1}}) a \qty{4}{Myr},
\qty{35}{m.Myr^{-1}} (\qty{0.035}{mm.yr^{-1}}) a \qty{25}{Myr}, \qty{25}{m.Myr^{-1}}
(\qty{0.025}{mm.yr^{-1}}) a \qty{49}{Myr}. (b) $4.2/0.0875 = 48$ veces. (c) $2.6/0.0875 = 30$
veces: la fase de carga hunde una isla mucho más deprisa que el enfriamiento de una placa joven.
\end{solution}

\begin{solution}{exo:b2:oceanic-islands:6}
(a) $\dot S = (3200/62.8)\,e^{-100/62.8} = 50.96\times0.2035 = \qty{10.37}{m}$ por millón de años,
\qty{0.0104}{m} por cada mil años: los \qty{0.01}{m} por cada mil años medidos. (b) $e^{-t/62.8}
= 10/50.96$, $t = 62.8\ln5.096 = \qty{102.3}{Myr}$. (c) Ley de la raíz: $2500 + 350\times5 =
\qty{4250}{m}$; ley exponencial: $6400 - 3200\times0.6716 = \qty{4251}{m}$; difieren en
\qty{1}{m}.
\end{solution}

\begin{solution}{exo:b2:oceanic-islands:7}
(a) $\qty{8.6}{cm.yr^{-1}} = \qty{86}{km.Myr^{-1}}$; $86\times33 = \qty{2838}{km}$. (b) Sí:
\qty{27.7}{Myr} es menos que los \qty{33}{Myr} tras los cuales se ahogan los arrecifes, así que su arrecife debería
seguir vivo. (c) $55\times1.609 = \qty{88.5}{km}$, o sea $88.5/86 = \qty{1.0}{Myr}$ más: unos
\qty{28.7}{Myr}. Es algo menos que \qty{33}{Myr}, pero esos \qty{33}{Myr} son a su vez una cifra
aproximada para toda la cadena, y el crecimiento depende también de la temperatura del agua en cada
\omterm{def:b2:oceanic-islands:reefs}{atolón}: Kure, el último \omterm{def:b2:oceanic-islands:reefs}{atolón} vivo, está en efecto cerca del \omterm{def:b2:oceanic-islands:darwin-point}{punto de Darwin}.
\end{solution}

\begin{solution}{exo:b2:oceanic-islands:8}
(a) Caja: $30\times17\times2 = \qty{1020}{km^3}$; pirámide: un tercio, \qty{340}{km^3}. (b) Un solo bloque
es la décima parte (caja) o la trigésima parte (pirámide) de todo el volcán Kilauea. (c) $23\,000\times
0.1 = \qty{2300}{km^3}$, una cuarta parte del Kilauea; una estimación que se sostiene o cae con el espesor
supuesto.
\end{solution}

\begin{solution}{exo:b2:oceanic-islands:9}
(a) $1000/2.6 = \qty{385}{kyr}$. (b) Sí: \qty{2.6}{m} por cada mil años son
\qty{2.6}{mm.yr^{-1}}, menos que \qty{4.2}{mm.yr^{-1}}, con un margen de \qty{1.6}{mm.yr^{-1}}.
(c) Subida relativa $12 + 2.6 = \qty{14.6}{mm.yr^{-1}}$; el arrecife se queda atrás $14.6 - 4.2 =
\qty{10.4}{mm.yr^{-1}}$ durante $16.5 - 8.2 = \qty{8.3}{kyr}$: su techo acaba a $10.4\times8.3 =
\qty{86}{m}$. (d) $(12 - 4.2)\times8.3 = \qty{65}{m}$. En ambos casos queda por encima del límite de crecimiento
más profundo, pero en aguas mucho más oscuras: que se recupere depende de cómo se ralentiza su crecimiento con la
profundidad.
\end{solution}

\begin{solution}{exo:b2:oceanic-islands:10}
(a) $S(t_1) = C(\sqrt{t_1 + \Delta} - \sqrt{t_1})$ tiene por derivada $\frac C2(1/\sqrt{t_1 +
\Delta} - 1/\sqrt{t_1}) < 0$: decrece con $t_1$. (b) $350\sqrt{20} = \qty{1565}{m}$; $350\times
(\sqrt{70} - \sqrt{50}) = 350\times1.2955 = \qty{453}{m}$. (c) No: $\sqrt t$ es finita en $0$, de modo que
el hundimiento en cualquier intervalo es finito, \qty{1565}{m} en (b). Solo la velocidad en el instante $t = 0$
es infinita, donde la ley ya no pretende describir el fondo en detalle.
\end{solution}

\begin{solution}{exo:b2:oceanic-islands:11}
(a) $30\times0.1 = \qty{3}{m}$. (b) Unos \qty{120}{m} si su techo estaba al nivel del mar en el
nivel alto: el \omterm{def:b2:oceanic-islands:reefs}{atolón} se convierte en una isla caliza, la lluvia disuelve el aragonito de los corales y la
roca recristaliza en calcita, con cuevas. (c) $120/3 = 40$ ciclos, \qty{4}{Myr}. (d) Las zonas lixiviadas
son esas épocas de exposición, selladas bajo el arrecife que volvió a crecer cuando regresó el mar.
A lo largo de muchos ciclos el mar vuelve a niveles altos parecidos, así que $\Delta\eta$ es pequeño, el techo del arrecife
vuelve a crecer hasta el nivel del mar en cada nivel alto, y el espesor del casquete lo fija el hundimiento, como
dice la proposición.
\end{solution}

\begin{solution}{exo:b2:oceanic-islands:12}
(a) $3200\times(e^{-50/62.8} - e^{-120/62.8}) = 3200\times(0.4510 - 0.1480) = \qty{970}{m}$.
(b) $150 + 970 = \qty{1120}{m}$. (c) $700\times1.8288 = \qty{1280}{m}$, \qty{160}{m} más profunda que la
estimación. (d) El \omterm{def:b2:oceanic-islands:seamount}{guyot} puede haberse ahogado antes, sobre una placa más joven, y haberse hundido más; un volcán
vecino puede haberlo cargado; su cima puede haber sido arrasada más abajo por el oleaje, o estar cubierta de sedimento; la
placa bajo él puede tener otra edad; y el nivel del mar del momento del ahogamiento no era
el actual.
\end{solution}

\begin{solution}{pb:b2:oceanic-islands:1}
\textbf{1.} $2500 + 350\sqrt{10} = 2500 + 350\times3.1623 = \qty{3607}{m}$.
\textbf{2.} $3607 + 1000 = \qty{4607}{m}$.
\textbf{3.} $V = \pi r^2h/3 = \pi\times50^2\times4.607/3 = \qty{12061}{km^3}$.
\textbf{4.} $0.85\times12\,061 = \qty{10252}{km^3}$ en \qty{1}{Myr}: \qty{0.0103}{km^3} al año,
diez veces menos que el Kilauea hoy.
\textbf{5.} $95\,600/12\,061 = 7.9$: unos ocho.
\textbf{6.} Lava toleítica, de la \omterm{def:b2:oceanic-islands:stages}{etapa de escudo}; las lavas alcalinas de la etapa de postescudo lo
recubren.
\textbf{7.} $\qty{1}{km}/\qty{1}{Myr} = \qty{1}{m}$ por cada mil años, 2.6 veces más lento que frente a
la isla de Hawái.
\textbf{8.} Sí: $4.2 - 1.0 = \qty{3.2}{mm.yr^{-1}}$ de sobra.
\textbf{9.} $12 + 1 - 4.2 = \qty{8.8}{mm}$ al año.
\textbf{10.} $150/8.8 = \qty{17}{kyr}$. Una terraza arrecifal ahogada en el flanco, un escalón del tipo de los
hallados frente a la isla de Hawái.
\textbf{11.} $350\times(\sqrt{22} - \sqrt{10}) = 350\times(4.6904 - 3.1623) = \qty{535}{m}$.
\textbf{12.} $1000 - 535 = \qty{465}{m}$, a $465/12 = \qty{38.8}{m}$ por millón de años.
\textbf{13.} $350/(2\sqrt{10}) = \qty{55.3}{m.Myr^{-1}}$ al principio; $350/(2\sqrt{22}) =
\qty{37.3}{m.Myr^{-1}}$ al final.
\textbf{14.} $2500 + 350\times4.6904 = \qty{4142}{m}$.
\textbf{15.} Las grandes islas hawaianas se elevan hasta unas cuatro veces más alto, de modo que hay mucha más
roca sobre el nivel del mar que eliminar; su \omterm{def:b2:oceanic-islands:stages}{etapa de escudo} duró además más tiempo.
\textbf{16.} El oleaje, que hace retroceder la costa y talla una plataforma de abrasión; los deslizamientos y los \omterm{def:b2:oceanic-islands:collapse}{colapsos
de flanco}.
\textbf{17.} Un \omterm{def:b2:oceanic-islands:reefs}{arrecife barrera}, con una laguna entre el arrecife y lo que queda de la costa volcánica.
\textbf{18.} $350\times(\sqrt{70} - \sqrt{22}) = 350\times(8.3666 - 4.6904) = \qty{1287}{m}$.
\textbf{19.} $1287/48 = \qty{26.8}{m}$ por millón de años, \qty{0.027}{mm} al año: $4.2/0.0268 =
157$ veces menos de lo que podría crecer el arrecife.
\textbf{20.} \qty{20}{m} más grueso que la caliza sobre el basalto bajo la isla Parry
(\qty{1267}{m}), y \qty{118}{m} más delgado que bajo Elugelab (\qty{1405}{m}): del mismo orden.
\textbf{21.} $2500 + 350\times8.3666 = \qty{5428}{m}$.
\textbf{22.} $3200\times(e^{-70/62.8} - e^{-120/62.8}) = 3200\times(0.3280 - 0.1480) =
\qty{576}{m}$.
\textbf{23.} $150 + 576 = \qty{726}{m}$.
\textbf{24.} El casquete arrecifal del \omterm{def:b2:oceanic-islands:reefs}{atolón} de Petrel tiene \textbf{\qty{1287}{m}} de espesor, unos \qty{1.3}{km}
de caliza arrecifal, como bajo Enewetak.
\end{solution}
